\chapter{The Foreign-Function Bridge}
\label{ch:ffi}

The C{--} kernel (Chapter~\ref{ch:kernel}) is fast but narrow; Lean~4
(Chapter~\ref{ch:lean}) is general but must trust external code to use it. The
FFI bridge reconciles them: it lets Lean call the kernel and, crucially,
\emph{cross-validates} the two. This chapter details the bridge.

\section{Architecture}
\label{sec:ffi-arch}

The bridge (\texttt{mathlib5-ffi-bridge/}) has three parts:
\begin{enumerate}
  \item \textbf{C side} (\texttt{C/src/bridge.c}, \texttt{C/include/bridge.h})
        --- a pure compute context with a rewrite table.
  \item \textbf{Lean side} (\texttt{Lean/Bridge.lean}, \texttt{Lean/TestBridge.lean})
        --- the FFI declarations and round-trip tests.
  \item \textbf{Theorem calculator} (\texttt{Lean/TheoremCalc/Main.lean})
        --- the cross-validation driver.
\end{enumerate}

\section{The C Context}
\label{sec:ffi-c}

\texttt{bridge.c} defines a \texttt{BridgeContext} with a config, a PRNG
state, a 256-entry rewrite table, and a 128-entry verification cache:
\begin{lstlisting}[language=C,caption={Bridge context (bridge.c).},label=lst:ffi-c]
struct BridgeContext {
    uint64_t config;
    int64_t  state;
    int64_t  rewrite_table[256];  // Pattern -> closed-form mapping
    int32_t  verify_cache[128];   // Theorem verification cache
};
\end{lstlisting}
The rewrite table is pre-populated with magic cookies for the three summation
patterns (SumSquares, SumLinear, SumCubes). \texttt{bridge\_rewrite} dispatches
on the pattern id to the closed form:
\begin{lstlisting}[language=C,caption={Closed-form dispatch (bridge.c).},label=lst:ffi-dispatch]
case 1: { int64_t n = input;
          return n * (n + 1) * (2 * n + 1) / 6; }   // SumSquares
case 2: { int64_t n = input;
          return n * (n + 1) / 2; }                  // SumLinear
case 3: { int64_t n = input; int64_t s = n*(n+1)/2;
          return s * s; }                            // SumCubes
\end{lstlisting}
The context is wiped on free (\texttt{bridge\_free}) --- sensitive state is
zeroed before deallocation, a small but deliberate hygiene measure.

\section{The Lean FFI Declarations}
\label{sec:ffi-lean}

\texttt{Bridge.lean} declares the foreign functions so Lean can call into the
compiled C. The context type and the three core calls
(\texttt{init}, \texttt{rewrite}, \texttt{free}) are surfaced as a
\texttt{BridgeCtx} wrapper used by the calculator.

\section{Cross-Validation: The Core Idea}
\label{sec:ffi-cross}

The theorem calculator computes each closed form two ways:
\begin{enumerate}
  \item \texttt{computeFFI} --- via the C kernel (fast path).
  \item \texttt{computeLean} --- independently via Lean's
        \texttt{patternRegistry} (slow but trusted path).
\end{enumerate}
It then asserts equality in \texttt{verifyConsistency}:
\begin{lstlisting}[language=Lean4,caption={Consistency check (TheoremCalc/Main.lean).},label=lst:ffi-consistency]
def verifyConsistency (ctx : BridgeCtx) (pid : Int64) (n : Int64) : IO Bool := do
  let ffiResult <- computeFFI ctx pid n
  let leanResult := computeLean pid n
  match leanResult with
  | some lr => pure (ffiResult == lr)
  | none => pure false
\end{lstlisting}

\section{The Main Driver}
\label{sec:ffi-main}

\texttt{Main.lean} registers the patterns, verifies FFI$\leftrightarrow$Lean
consistency at $n=10$, benchmarks \texttt{SumSquares} for $n=1..1000$, and
prints sample results. The benchmark loop:
\begin{lstlisting}[language=Lean4,caption={Benchmark loop (TheoremCalc/Main.lean).},label=lst:ffi-bench]
let start <- IO.monoMsNow
let mut result := (0 : Int64)
for i in List.range 1000 do
  let n := Int64.ofNat (i + 1)
  result <- computeFFI ctx 1 n
let elapsed <- IO.monoMsNow
\end{lstlisting}
A passing run ends with \texttt{=== All verifications passed ===}, which the
CI gate (\cref{ch:build}) greps for.

\section{Why Cross-Validation Matters}
\label{sec:ffi-why}

A single implementation can be wrong and confident. Two independent
implementations that agree on every tested input are vastly more trustworthy.
If they ever disagree, the discrepancy is a receipt, not a silent error. This
is the same principle as the P/NP swarm (hard to find, easy to check)
applied to arithmetic: Lean is the cheap, trustworthy checker; the kernel is
the fast solver.

\section{Build Orchestration}
\label{sec:ffi-build}

\texttt{scripts/nats\_build\_orchestrator.sh} drives the multi-language build
(C $\ensuremath{rightarrow}$ shared object, Lean $\ensuremath{rightarrow}$ \texttt{lake build},
orchestration). The \texttt{lean-toolchain} pins the Lean version so the FFI
ABI is reproducible --- a receipt (\cref{ch:worm}) records exactly which
toolchain produced a binary.

\section{Foreign Call Contract}
\label{sec:ffi-contract}

The C$\leftrightarrow$Lean boundary obeys a strict calling convention:
Lean passes a tagged value (integer or float) and receives a tagged result;
the C side never observes Lean's heap. The bridge checks the tag on entry,
so a value of the wrong type is rejected with a sealed error receipt rather
than causing undefined behavior. This is the only place where the two
memory models meet, and confining it to one audited module bounds the
trusted surface of the heterogeneous stack.
